\section{Repairing a deployed tokenizer}
\label{sec:method}

The repair has three parts: a one-class edit to the regex, new merges learned on top of the existing merge list, and an embedding initialisation that leaves the model's computation on old tokens unchanged.

\paragraph{Regex edit.} We admit \pM{} wherever the regex has a letter word class, and exclude it from the classes that would otherwise absorb a mark as a prefix or as punctuation. For the Llama-3/Qwen pattern this is two substitutions:
\begin{center}\small
\texttt{[\textasciicircum\textbackslash r\textbackslash n\textbackslash p\{L\}\textbackslash p\{N\}]?\textbackslash p\{L\}+}\\
$\downarrow$\\
\texttt{[\textasciicircum\textbackslash r\textbackslash n\textbackslash p\{L\}\textbackslash p\{M\}\textbackslash p\{N\}]?[\textbackslash p\{L\}\textbackslash p\{M\}]+}
\end{center}
and the same change in the punctuation branch \texttt{\textvisiblespace?[\textasciicircum\textbackslash s\textbackslash p\{L\}\textbackslash p\{N\}]+}. The optional leading class and the punctuation class must exclude \pM{}, so that a mark can only be matched by the word class and stays in one pre-token with the letters around it. Nothing else in the pattern changes, so digits, contractions, whitespace and punctuation are handled as before. We call the result the \emph{repaired regex}. Qwen3.5 independently shipped almost the same word branch (\S\ref{sec:audit};~\citealp{land2026smallbugs}).

\paragraph{Continued BPE, restricted to Devanagari.} Following \citet{purason-etal-2026-teaching}, we learn $K$ new merges on top of the existing ones. We pre-tokenize Nepali text with the arm's regex and segment every pre-token with the \emph{original} merges, which is how the deployed model sees it. Standard BPE over these segmentations, treating each base token as an indivisible symbol, then yields new merges, which we append after the existing ones. We discard any merge whose output contains no code point in U+0900--U+097F, together with merges that depend on it. At inference, BPE applies the lowest-ranked applicable merge first, so the original merges act before the new ones. The extended tokenizer reproduces the training procedure except in two cases: a new merge whose output string is already in the original vocabulary can let an original merge fire again afterwards, and under HuggingFace's \texttt{ignore\_merges} option (set for Llama-3) a pre-token that is itself a vocabulary entry is emitted directly. Neither case affects Proposition~\ref{prop:lossless}. We run the merge learning inside the HuggingFace Rust trainer by mapping each base token to a private-use code point. Learning $K{=}32{,}000$ merges on 1\,GB of Nepali text takes 1.7 to 3.4 minutes per tokenizer.

\RI{} and \RII{} use this identical procedure on identical text with identical $K$. The only difference is the regex used to pre-tokenize, both when learning merges and at inference.

\begin{proposition}[English invariance]
\label{prop:lossless}
Let $s$ be a string, after the tokenizer's normalizer, containing no code point of category \pM{} and no code point in U+0900--U+097F. Then the repaired tokenizer \RII{} and the original tokenizer produce identical token ids for $s$.
\end{proposition}
\begin{proof}
\emph{Pre-tokens.} The repaired regex differs from the original only in the three classes edited above, and each edited class differs from its original only on \pM{} characters. A regex match on $s$ consults only characters of $s$, and the order of the alternatives is unchanged. Since $s$ has no \pM{} character, both regexes produce the same pre-tokens.
\emph{Merges and vocabulary.} \RII{} differs from the original tokenizer only by appended merges and by the vocabulary entries they create, each of which contains a complete code point in U+0900--U+097F. A merge $(a,b)$ applies only where $a$ and $b$ are adjacent symbols of a pre-token, so its output is a contiguous substring of that pre-token's bytes. A vocabulary entry is emitted directly under \texttt{ignore\_merges} only if it equals the whole pre-token. Neither can happen for a pre-token of $s$, because UTF-8 is self-synchronising and $s$ contains no such code point. Each pre-token of $s$ is therefore encoded with original merges and original vocabulary entries only, exactly as by the original tokenizer.
\emph{Ids.} New ids are appended above the original ids, so all ids coincide.
\end{proof}

The proposition covers ASCII English and code, and every other script whose normalised text has no combining marks, including most Cyrillic, Han, Hangul and Ethiopic text. It does not cover Devanagari, even without marks: mark-free Devanagari words and Devanagari digits contain code points in U+0900--U+097F and change by design, since new merges apply to them. It also does not cover text with marks, such as Hindi or Latin written in decomposed form (NFD, e.g.\ \emph{caf\'e} with a combining acute). The Llama-3 tokenizer has no normalizer, so NFD input to a deployed Llama-3 \RII{} model falls outside the guarantee; Qwen3's tokenizer applies NFC itself, which composes most, but not all, combining marks. We normalise all our text to NFC. Empirically, \RI{} and \RII{} at $K{=}32{,}000$ give token ids identical to the original tokenizer on all 1{,}012 FLORES devtest sentences in English, Russian and Amharic, for both models. One Russian sentence contains a combining acute (U+0301) and so lies outside the proposition; its ids are nonetheless identical.

\paragraph{Embedding initialisation.} Each new token's constituents are the base tokens obtained by applying the original merges to the token's byte-level string; we do not decode it to text, which would turn tokens ending in a partial UTF-8 sequence into U+FFFD. Its input and output embedding rows are set to the mean of the constituent rows~\citep{gee-etal-2022-fast}. Original rows are untouched. Qwen3's embedding matrix has unused padding rows after its vocabulary, and the first new tokens reuse them. On any input whose ids are all original, the resized model therefore produces the same hidden states, and the same logits for every original token, as the base model. We check this numerically before training on an English probe sentence: the original-vocabulary logits are identical up to float error (maximum absolute difference $<10^{-4}$). The softmax denominator gains the new tokens' logits, so output probabilities are not identical, and continued pretraining is still needed.

\paragraph{Arms.} \RO{} (no extension) keeps the original tokenizer. \RI{} (extend) appends $K$ merges learned under the original letters-only regex: the standard recipe~\citep{purason-etal-2026-teaching,yamaguchi2026lowres}. \RII{} (repair + extend) switches to the repaired regex and appends $K$ merges learned under it. All arms then continue pretraining on the same bytes (\S\ref{sec:setup}).
